\documentclass[14pt,a4paper,final]{extreport}
\nonstopmode
\usepackage[T2A]{fontenc}
\usepackage[utf8]{inputenc}
\usepackage{cmap}
\DeclareUnicodeCharacter{0306}{~}
\usepackage[russian]{babel}
\usepackage[a4paper, top=2cm, bottom=2cm, left=2cm, right=1cm,marginparsep=.5cm]{geometry}
\usepackage{fancyhdr}

\usepackage[export]{adjustbox}
\usepackage{marginnote}
\usepackage{tabularray}
\usepackage{longtable}
\usepackage{ltablex}

\usepackage{enumitem}
\setlist[itemize]{leftmargin=\parindent,nosep}
\setlist[enumerate]{leftmargin=\parindent,nosep}
\renewcommand\labelitemi{--}

\usepackage[none]{hyphenat}
\usepackage{setspace}
\onehalfspacing 
\usepackage{float}
\usepackage{xpatch}
\usepackage{moreverb}

\usepackage{tocloft, titlesec, titletoc}
\usepackage{indentfirst}


\usepackage{listings}
\usepackage{xcolor}

% Настройка глобального стиля для ВСЕХ блоков кода
\lstset{
    basicstyle=\ttfamily\footnotesize, % ЕДИНЫЙ размер шрифта для всего кода (примерно 10pt)
    breaklines=true,                   % ВКЛЮЧАЕТ АВТОМАТИЧЕСКИЙ ПЕРЕНОС СТРОК!
    breakatwhitespace=false,           % Переносить строки даже посреди длинных выражений
    numbers=left,                      % Нумерация строк слева
    numberstyle=\tiny\color{gray},     % Стиль номеров строк
    stepnumber=1,                      % Нумеровать каждую строку
    frame=single,                      % Красивая рамка вокруг кода (по желанию, можно убрать)
    tabsize=4,                         % Размер табуляции
    showstringspaces=false,            % Не показывать уродливые символы пробелов в строках
    keepspaces=true                    % Сохранять пробелы и отступы C++
}



\titleformat{\chapter}[block]{\normalfont\bfseries\large}{\thechapter}{1ex}{\centering\MakeUppercase}
\titlespacing*{\chapter}{0pt}{-\baselineskip}{2\baselineskip}

\titleformat{\section}[block]{\normalfont}{\thesection}{1ex}{\normalfont}
\titlespacing{\section}{1.25cm}{\baselineskip}{\baselineskip}

%Оформление подразделов
\titleformat{\subsection}[block]{\normalfont}{\thesubsection}{1ex}{\normalfont}
\titlespacing{\subsection}{1.25cm}{\baselineskip}{\baselineskip}

%Оформление под-подразделов
\titleformat{\subsubsection}[block]{\normalfont}{\thesubsection}{1ex}{\normalfont}
\titlespacing{\subsubsection}{1.25cm}{0pt}{\baselineskip}

\usepackage{titlesec}

% Настройка глубины нумерации заголовков
\setcounter{secnumdepth}{3}
\setcounter{tocdepth}{3}

% Настройка внешнего вида заголовков с помощью titlesec
\titleformat{\subsubsection}[hang]{\normalfont\normalsize}{\thesubsubsection}{0.5em}{}

%\titlespacing{\subsubsection}{35pt}{\parskip}{\parskip}

%\setcounter{tocdepth}{4}%
%\setcounter{secnumdepth}{4}% Number \subsubsection
\renewcommand{\cftdotsep}{2.75}
\renewcommand{\cftchapfont}{\normalfont}
\renewcommand{\cfttoctitlefont}{\hfill\normalfont\large\bfseries}
\renewcommand{\cftaftertoctitle}{\hfill}
\renewcommand{\cftchapleader}{\cftdotfill{\cftdotsep}} % for chapters
\renewcommand{\cftsecleader}{\cftdotfill{\cftdotsep}}  % for sections
\renewcommand{\cftsubsecleader}{\cftdotfill{\cftdotsep}} % for subsections
\renewcommand{\cftsubsubsecleader}{\cftdotfill{\cftdotsep}} % for subsubsections
\setlength{\cftbeforetoctitleskip}{-.6\baselineskip}
\setlength{\cftaftertoctitleskip}{-\baselineskip}
\setlength{\cftbeforechapskip}{\baselineskip}
\renewcommand{\cftchappagefont}{}

\renewcommand{\contentsname}{СОДЕРЖАНИЕ\vspace{\baselineskip}}% (ГОСТ Р 7.0.11-2011, 4)
\addto{\captionsrussian}{\renewcommand{\bibname}{СПИСОК ИСПОЛЬЗОВАННОЙ ЛИТЕРАТУРЫ\bigskip}}

\makeatletter
\renewcommand\@biblabel[1]{#1.\ }
\makeatother

\let\mtcontentsname\contentsname\addto\captionsrussian{\typeout{foo}\renewcommand{\contentsname}{\MakeUppercase\mtcontentsname}}

\makeatletter
\patchcmd{\chapter}{\if@openright\cleardoublepage\else\clearpage\fi}{\vspace{3\baselineskip}}{}{}
\makeatother

\newcommand*\rot[1]{\rotatebox{90}{#1}}
\newlength{\rs}
\setlength{\rs}{.5mm}

\makeatletter
\AddEnumerateCounter{\asbuk}{\russian@alph}{щ}
\makeatother

\fancypagestyle{firststyle}
{
    \fancyhf{} % sets both header and footer to nothing
    \renewcommand{\headrulewidth}{0pt}
    \fancyfoot[C]{
    
        {\vspace{-\baselineskip}\normalfont\textbf{Москва 2026}}
    } 
}

\begin{document}
\pagenumbering{gobble}
\raggedbottom

%для изменения
\newcommand{\fullname}{Нагороднюк Алиса Дмитриевна}
\newcommand{\shname}{Нагороднюк А. Д.}
\newcommand{\grp}{М3О-225БВ-24}
\newcommand{\variant}{9} %номер варианта (если многозначный - уберите 0 в \cipher)
\newcommand{\fieldA}{\( \vec a = (z^3) \vec i + (x^3) \vec j + (y^3) \vec k \)}
\newcommand{\fieldB}{\( \vec b = (y^3z^4 + 2xy) \vec i + (3xy^2z^4 + x^2) \vec j + (4xy^3z^3 + 2z) \vec k\)}

\newcommand{\cipher}{643.02066606.000\variant}
\newcommand{\cipherlri}{643.02066606.\\000\variant}
\newcommand{\docstamp}{\cipher-01 12 01}
\newcommand{\doctitle}{Программа расчета\\параметров векторных полей\\(Теория поля)}
\newcommand{\doctitleoneline}{Программа расчета параметров векторных полей (Теория поля)}
\newcommand{\programtitle}{PR\_BLACK}
\newcommand{\docname}{Текст программы}
\input{lu.tex}
\pagebreak

\input{title.tex}
\pagebreak

\parindent=1.25cm
\sloppy

\pagenumbering{arabic}
\setcounter{page}{2}

\pagestyle{fancy}

\fancypagestyle{plain}

\fancyhf{} % clear all header and footer fields

\fancyhead[C]{\bfseries\thepage\\\docstamp}

\renewcommand{\headrulewidth}{0pt}
\pagebreak

\chapter{ТЕКСТ ПРОГРАММЫ \programtitle\ НА ИСХОДНОМ ЯЗЫКЕ}

\section{Файл заголовков Resource.h}
\footnotesize
\begin{listing}{1}
/*************************************************************************
* КАФЕДРА № 304 | 2 КУРС | УЧЕБНАЯ ПРАКТИКА                             *
*------------------------------------------------------------------------*
* Project Name  : Программа расчета параметров векторных полей           *
* (Теория поля)                                                          *
* File Name     : Resource.h                                             *
* Author(s)     : Варфоломеев К.В., Нагороднюк А.Д.,                     *
*                 Великоцкая А.С.                                        *
* Created       : 02/07/2026                                             *
* Last Modified : 09/07/2026                                             *
* Version       : 1.0                                                    *
* Copyright     : (c) 2026. All rights reserved.                         *
* License       : Educational use only                                   *
*------------------------------------------------------------------------*
* НАЗНАЧЕНИЕ ФАЙЛА:                                                      *
* Содержит макроопределения идентификаторов графических ресурсов и       *
* элементов управления пользовательского интерфейса Win32.               *
*                                                                        *
* ОПИСАНИЕ ФУНКЦИОНИРОВАНИЯ:                                             *
* Файл Resource.h автоматически генерируется средой разработки           *
* Microsoft Visual Studio и дополняется разработчиком для связывания     *
* визуальных компонентов (кнопок, диалоговых окон, иконок) с их          *
* числовыми идентификаторами в таблице ресурсов приложения.              *
*                                                                        *
* В ДАННОМ ФАЙЛЕ ОПРЕДЕЛЕНЫ:                                             *
* 1. Системные идентификаторы приложения (IDS_APP_TITLE,                 *
* IDR_MAINFRAME, IDI_VECTORFIELD, IDI_SMALL).                            *
* 2. Идентификаторы диалоговых окон меню и подсистемы "О программе"      *
* (IDD_VECTORFIELD_DIALOG, IDD_ABOUTBOX).                                *
* 3. Уникальные идентификаторы управляющих кнопок на панели              *
* инструментов, обеспечивающие обработку пользовательских событий        *
* в оконной процедуре:                                                   *
* - ID_BTN_IMPORT : Импорт данных из внешнего файла .txt                 *
* - ID_BTN_SHOW   : Отображение матрицы полиномов                        *
* - ID_BTN_CALC   : Запуск процедур математического анализа полей        *
* - ID_BTN_EXPORT : Экспорт сформированного отчета на диск               *
*************************************************************************/

//{{NO_DEPENDENCIES}}
// Включаемый файл, созданный в Microsoft Visual C++.
// Используется VectorField.rc
#define IDS_APP_TITLE           103
#define IDR_MAINFRAME           128
#define IDD_VECTORFIELD_DIALOG  102
#define IDD_ABOUTBOX            103
#define IDM_ABOUT               104
#define IDM_EXIT                105
#define IDI_VECTORFIELD         107
#define IDI_SMALL               108
#define IDC_VECTORFIELD         109
#define IDC_MYICON              2
#ifndef IDC_STATIC
#define IDC_STATIC              -1
#endif
// Кнопки интерфейса
#define ID_BTN_IMPORT           1001
#define ID_BTN_SHOW             1002
#define ID_BTN_CALC             1003
#define ID_BTN_VIS_A            1004
#define ID_BTN_VIS_B            1005
#define ID_BTN_EXPORT           1006
// Следующие стандартные значения для новых объектов
#ifdef APSTUDIO_INVOKED
#ifndef APSTUDIO_READONLY_SYMBOLS
#define _APS_NO_MFC             130
#define _APS_NEXT_RESOURCE_VALUE 129
#define _APS_NEXT_COMMAND_VALUE 32771
#define _APS_NEXT_CONTROL_VALUE 1007
#define _APS_NEXT_SYMED_VALUE   110
#endif
#endif
\end{listing}
\normalsize

\pagebreak

\section{Файл заголовков VectorField.h}
\footnotesize
\begin{listing}{1}
/*************************************************************************
* КАФЕДРА № 304 | 2 КУРС | УЧЕБНАЯ ПРАКТИКА                             *
*------------------------------------------------------------------------*
* Project Name  : Программа расчета параметров векторных полей           *
* (Теория поля)                                                          *
* File Name     : VectorField.h                                          *
* Author(s)     : Варфоломеев К.В., Нагороднюк А.Д.,                     *
*                 Великоцкая А.С.                                        *
* Created       : 02/07/2026                                             *
* Last Modified : 09/07/2026                                             *
* Version       : 1.0                                                    *
* Copyright     : (c) 2026. All rights reserved.                         *
* License       : Educational use only                                   *
*------------------------------------------------------------------------*
* НАЗНАЧЕНИЕ ФАЙЛА:                                                      *
* Содержит объявления базовых структур данных трехмерного пространства,  *
* класса символьных полиномов и статических методов аналитического       *
* ядра для вычислений в теории поля.                                     *
*                                                                        *
* ОПИСАНИЕ ИНТЕРФЕЙСОВ И ТИПОВ ДАННЫХ:                                   *
* В заголовочном файле сосредоточены все типы данных, на которых         *
* базируется вычислительная логика программы:                            *
*                                                                        *
* 1. Структура Vector3D:                                                 *
* Описывает трехмерный геометрический вектор с вещественными             *
* координатами (x, y, z).                                                *
*                                                                        *
* 2. Структура Monomial:                                                 *
* Описывает структуру отдельного монома, содержащего вещественный        *
* коэффициент "coef" и целочисленные степени трех пространственных       *
* переменных (px, py, pz).                                               *
*                                                                        *
* 3. Класс Polynomial:                                                   *
* Инкапсулирует логику символьной работы со многочленами (сложение,      *
* зануление переменных при расчете потенциалов, символьное               *
* интегрирование, численное вычисление значения в точке).                *
*                                                                        *
* 4. Класс FieldMathCalculations:                                        *
* Математическое ядро программы, выполняющее синтаксический анализ       *
* строковых формул, расчет дивергенции, ротора, потока (через теорему    *
* Гаусса), циркуляции (через теорему Стокса), а также нахождение         *
* потенциала векторного поля.                                            *
*************************************************************************/
#pragma once
#include <string>
#include <vector>
#include <functional>
#include "resource.h"
using namespace std;
// =======================================================================
// ОБЪЯВЛЕНИЯ СТРУКТУР И КЛАССОВ
// =======================================================================
struct Vector3D {double x, y, z;}; struct Monomial {double coef; int px, py, pz;};
class Polynomial {
public:
    vector<Monomial> terms; Polynomial(initializer_list<Monomial> init);
    Polynomial(); Polynomial SetVariableToZero(bool zeroX, bool zeroY,
    bool zeroZ) const;
    Polynomial IntegrateSymbolically(int varIndex) const;
    void Add(const Polynomial& other); wstring ToWString() const;
    double Evaluate(double x, double y, double z) const;
};
class FieldMathCalculations {
public: typedef function<double(double, double, double)> FUNC3D;
    typedef function<Vector3D(double, double, double)> FUNC3D_VEC;
    static Polynomial fieldA_X, fieldA_Y, fieldA_Z;
    static Polynomial fieldB_X, fieldB_Y, fieldB_Z;
    static void ParsePolynomials(const wstring& a_str, const wstring& b_str);
    static Vector3D EvaluateFieldA(double x, double y, double z);
    static Polynomial GetFieldBX(); static Polynomial GetFieldBY();
    static Polynomial GetFieldBZ();
    static Vector3D EvaluateFieldB(double x, double y, double z);
    static double Calc_PartialDerivative(FUNC3D f, int variable, 
                                       double x, double y, double z);
    static double IntegrateQuarterDisk(function<double(double, double)> func, 
                                       double radius = 1.0, int steps = 50);
    static double Calc_Divergence(FUNC3D_VEC field, double x, double y, double z);
    static Vector3D Calc_Rotor(FUNC3D_VEC field, double x, double y, double z);
    static wstring DetermineFieldClass(FUNC3D_VEC field);
    static double Calc_Flux_GaussOstrogradsky();
    static double Calc_Flux_Direct(); static double Calc_Circulation_Stokes();
    static double Calc_Circulation_Direct(); static Polynomial GetAnalyticalPotential();
    static bool VerifyPotentialCorrectness(); static bool CheckRotorBZero();
private: static Polynomial ParsePoly(wstring expr);
    static void ExtractComponents(wstring line, Polynomial& P, 
                                  Polynomial& Q, Polynomial& R);
    static void ReplaceAll(wstring& str, const wstring& from, const wstring& to);
};
\end{listing}
\normalsize

\pagebreak

\section{Метод Polynomial::SetVariableToZero}
\footnotesize
\begin{listing}{1}
/*************************************************************************
* КАФЕДРА № 304 | 2 КУРС | УЧЕБНАЯ ПРАКТИКА                              *
*------------------------------------------------------------------------*
* Проект        : Программа расчета параметров векторных полей           *
* Файл          : VectorField.cpp                                        *
* Создан        : 02.07.2026                                             *
*                                                                        *
* РАСПРЕДЕЛЕНИЕ ОТВЕТСТВЕННОСТИ (ФУНКЦИОНАЛЬНЫЕ МОДУЛИ):                 *
* - Великоцкая А.С. : Модули 1-20 (Математическое ядро)                  *
* - Варфоломеев К.В. : Модули 20-35 (Вычисления и интерфейс)             *
* - Нагороднюк А.Д. : Модули 35-49 (Интерфейс и визуализация)            *
*------------------------------------------------------------------------*
* НАЗНАЧЕНИЕ ФАЙЛА:                                                      *
* Реализация логики работы системы: описание полей, алгоритмы анализа,   *
* обработка данных и управление интерфейсом (WinAPI).                    *
*************************************************************************/
/*************************************************************************
* Метод         : Polynomial::SetVariableToZero                          *
*------------------------------------------------------------------------*
* Назначение    : Символьное зануление переменных в полиноме             *
* Входные       : zeroX, zeroY, zeroZ - флаги зануления осей x, y, z     *
* Выходные      : Новый объект Polynomial с отфильтрованными мономами    *
*------------------------------------------------------------------------*
* ОПИСАНИЕ ФУНКЦИОНИРОВАНИЯ:                                             *
* Метод осуществляет символьное зануление выбранных переменных в         *
* структуре полинома. Это необходимо для подготовки к аналитическому     *
* интегрированию (потенциала поля). Метод фильтрует мономы, отбрасывая   *
* те, у которых степени зануляемых переменных больше нуля.               *
*************************************************************************/
Polynomial Polynomial::SetVariableToZero(bool zeroX, bool zeroY, 
                                         bool zeroZ) const {
    Polynomial result;
    for (const auto& term : terms) {
        if ((zeroX && term.px > 0) || (zeroY && term.py > 0) || 
            (zeroZ && term.pz > 0)) continue;
        result.terms.push_back(term);
    }
    return result;
}
\end{listing}
\normalsize

\pagebreak

\section{Метод Polynomial::IntegrateSymbolically}
\footnotesize
\begin{listing}{1}
/*************************************************************************
* Метод         : Polynomial::IntegrateSymbolically                      *
*------------------------------------------------------------------------*
* Назначение    : Символьное интегрирование полинома по оси              *
* Входные       : varIndex - индекс переменной (0 -> x, 1 -> y, 2 -> z)  *
* Выходные      : Новый объект Polynomial (первообразная)                *
*------------------------------------------------------------------------*
* ОПИСАНИЕ ФУНКЦИОНИРОВАНИЯ:                                             *
* Выполняет аналитическое интегрирование полинома по выбранной оси.      *
* Метод обходит каждый моном и применяет правило интегрирования для      *
* степенной функции. Результат — новая первообразная функция,            *
* используемая для точного расчета потенциала.                           *
*************************************************************************/
Polynomial Polynomial::IntegrateSymbolically(int varIndex) const {
    Polynomial result;
    for (const auto& term : terms) {
        Monomial m = term;
        if (varIndex == 0) { m.px += 1; m.coef /= (double)m.px; }
        else if (varIndex == 1) { m.py += 1; m.coef /= (double)m.py; }
        else if (varIndex == 2) { m.pz += 1; m.coef /= (double)m.pz; }
        result.terms.push_back(m);
    }
    return result;
}
\end{listing}
\normalsize

\pagebreak

\section{Метод Polynomial::Add}
\footnotesize
\begin{listing}{1}
/*************************************************************************
* Метод         : Polynomial::Add                                        *
*------------------------------------------------------------------------*
* Назначение    : Алгебраическое сложение полиномов                      *
* Входные       : other - ссылка на складываемый полином                 *
* Выходные      : Нет (модифицирует текущий объект)                      *
*------------------------------------------------------------------------*
* ОПИСАНИЕ ФУНКЦИОНИРОВАНИЯ:                                             *
* Реализует сложение многочленов с приведением подобных членов.           *
* Алгоритм: 1. Поиск подобных слагаемых по набору степеней (px, py, pz).  *
* 2. Очистка от нулевых членов с коэффициентом < 1e-6 (микропогрешности). *
*************************************************************************/
void Polynomial::Add(const Polynomial& other) {
    for (const auto& otherTerm : other.terms) {
        bool found = false;
        for (auto& myTerm : terms) {
            if (myTerm.px == otherTerm.px && myTerm.py == otherTerm.py && 
                myTerm.pz == otherTerm.pz) {
                myTerm.coef += otherTerm.coef; found = true; break;
            }
        }
        if (!found) terms.push_back(otherTerm);
    }
    terms.erase(remove_if(terms.begin(), terms.end(), 
        [](const Monomial& m) { return abs(m.coef) < 1e-6; }), 
        terms.end());
}
\end{listing}
\normalsize

\pagebreak

\section{Метод Polynomial::Evaluate}
\footnotesize
\begin{listing}{1}
/*************************************************************************
* Метод         : Polynomial::Evaluate                                   *
*------------------------------------------------------------------------*
* Назначение    : Нахождение численного значения полинома в точке        *
* Входные       : x, y, z - координаты расчетной точки                   *
* Выходные      : Вещественное число double - значение в точке           *
*------------------------------------------------------------------------*
* ОПИСАНИЕ ФУНКЦИОНИРОВАНИЯ:                                             *
* Производит численное вычисление значения многочлена в заданной точке   *
* пространства. Математическая модель: вычисление суммы произведений     *
* коэффициентов на переменные в степенях. Используется функция pow       *
* из заголовочного файла <cmath>.                                        *
*************************************************************************/
double Polynomial::Evaluate(double x, double y, double z) const {
    double sum = 0.0;
    for (const auto& t : terms) 
        sum += t.coef * pow(x, t.px) * pow(y, t.py) * pow(z, t.pz);
    return sum;
}
\end{listing}
\normalsize

\pagebreak%
\section{Метод Polynomial::ToWString}
\footnotesize
\begin{listing}{1}
/*************************************************************************
* Метод         : Polynomial::ToWString                                  *
*------------------------------------------------------------------------*
* Назначение    : Форматирование полинома в строковое представление      *
* Входные       : Нет                                                    *
* Выходные      : wstring - строка широких символов (формат мономов)     *
*------------------------------------------------------------------------*
* ОПИСАНИЕ ФУНКЦИОНИРОВАНИЯ:                                             *
* Преобразует внутреннее представление полинома в строку вида            *
* "ax^2 + by - cz^3" для вывода в графическом интерфейсе.                *
* *
* АЛГОРИТМ:                                                              *
* 1. Если список "terms" пуст, возвращается строковый ноль ("0").         *
* 2. Итеративно обходится каждый моном. Определяются знаки слагаемых.    *
* 3. Форматируется вещественный коэффициент (std::setprecision(4)).      *
* 4. Единичные коэффициенты (при наличии переменных) опускаются.         *
* 5. Дописываются переменные x, y, z со степенями (если > 1).            *
*************************************************************************/
wstring Polynomial::ToWString() const {
    if (terms.empty()) return L"0";
    wstringstream ss; bool first = true;
    for (const auto& term : terms) {
        if (term.coef > 0 && !first) ss << L" + ";
        if (term.coef < 0) ss << (first ? L"-" : L" - ");
        double absCoef = abs(term.coef);
        if (absCoef != 1.0 || (term.px == 0 && term.py == 0 && 
            term.pz == 0)) ss << setprecision(4) << absCoef;   
        if (term.px > 0) { ss << L"x"; if (term.px > 1) ss << L"^" << term.px; }
        if (term.py > 0) { ss << L"y"; if (term.py > 1) ss << L"^" << term.py; }
        if (term.pz > 0) { ss << L"z"; if (term.pz > 1) ss << L"^" << term.pz; }
        first = false;
    }
    return ss.str();
}
\end{listing}
\normalsize

\pagebreak

\section{Функция FieldMathCalculations::ReplaceAll}
\footnotesize
\begin{listing}{1}
/*************************************************************************
* Функция       : FieldMathCalculations::ReplaceAll                      *
*------------------------------------------------------------------------*
* Назначение    : Глобальная замена всех вхождений подстроки в строке    *
* Входные       : str - модифицируемая строка,                           *
* from - искомая подстрока, to - заменяющая              *
* Выходные      : Нет (модифицирует строку по ссылке)                    *
*------------------------------------------------------------------------*
* ОПИСАНИЕ ФУНКЦИОНИРОВАНИЯ:                                             *
* Осуществляет поиск и замену всех вхождений подстроки "from" на         *
* значение "to" внутри строки "str".                                     *
*************************************************************************/
void FieldMathCalculations::ReplaceAll(wstring& str, 
                                       const wstring& from, 
                                       const wstring& to) {
    if (from.empty()) return;
    size_t start_pos = 0;
    while ((start_pos = str.find(from, start_pos)) != wstring::npos) {
        str.replace(start_pos, from.length(), to);
        start_pos += to.length();
    }
}
\end{listing}
\normalsize

\pagebreak

\section{Функция FieldMathCalculations::ParsePoly}
\footnotesize
\begin{listing}{1}
/*************************************************************************
* Функция       : FieldMathCalculations::ParsePoly                       *
*------------------------------------------------------------------------*
* Назначение    : Синтаксический анализ строкового полинома              *
* Входные       : expr - анализируемая строка формулы                    *
* Выходные      : Объект Polynomial (структура мономов)                  *
*------------------------------------------------------------------------*
* ОПИСАНИЕ ФУНКЦИОНИРОВАНИЯ:                                             *
* Разбирает математическое выражение полинома от x, y, z.                *
* *
* АЛГОРИТМ:                                                              *
* 1. Очистка строки от пробелов/скобок. Замена "-" на "+-".              *
* 2. Разделение строки на токены по "+".                                 *
* 3. Извлечение коэффициента и степеней переменных (x, y, z).            *
*************************************************************************/
Polynomial FieldMathCalculations::ParsePoly(wstring expr) {
    Polynomial poly; wstring clean;
    for (wchar_t c : expr) {
        if (c == L' ' || c == L'(' || c == L')' || c == L'*') continue;
        if (c == L'-') { clean += L"+-"; continue; } clean += c; }
    wstringstream ss(clean); wstring token;
    while (getline(ss, token, L'+')) {
        if (token.empty()) continue;
        Monomial m = { 1.0, 0, 0, 0 }; 
        size_t i = (token[0] == L'-') ? 1 : 0;
        wstring coef = (token[0] == L'-') ? L"-" : L"";
        while (i < token.length() && (iswdigit(token[i]) || 
               token[i] == L'.')) coef += token[i++]; 
        m.coef = (coef == L"-" || coef.empty()) ? 
                 ((coef == L"-") ? -1.0 : 1.0) : stod(coef);    
        auto pV = [&](wchar_t v, int& p) { size_t pos = token.find(v); 
            if (pos != wstring::npos) { p = 1; 
                if (pos + 1 < token.length() && token[pos + 1] == L'^') 
                    p = stoi(token.substr(pos + 2)); } 
        };
        pV(L'x', m.px); pV(L'y', m.py); pV(L'z', m.pz);
        poly.terms.push_back(m);}
    return poly;
}
\end{listing}
\normalsize

\pagebreak

\section{Функция FieldMathCalculations::ExtractComponents}
\footnotesize
\begin{listing}{1}
/*************************************************************************
* Функция       : FieldMathCalculations::ExtractComponents               *
*------------------------------------------------------------------------*
* Назначение    : Выделение координатных функций P, Q, R                 *
* Входные       : line - строка формулы, P, Q, R - ссылки на полиномы    *
* Выходные      : Нет                                                    *
*------------------------------------------------------------------------*
* ОПИСАНИЕ ФУНКЦИОНИРОВАНИЯ:                                             *
* Извлекает полиномы P, Q, R из строки векторного поля.                  *
* 1. Нормализация LaTeX-символики.                                       *
* 2. Разбор строки при наличии ортов (i, j, k) или через разделитель ";". *
*************************************************************************/
void FieldMathCalculations::ExtractComponents(wstring line, Polynomial& P, 
                                              Polynomial& Q, Polynomial& R) {
    ReplaceAll(line, L"\\overline{i}", L"i"); 
    ReplaceAll(line, L"\\overline{j}", L"j");
    ReplaceAll(line, L"\\overline{k}", L"k"); 
    ReplaceAll(line, L"\\overline{a}=", L""); 
    ReplaceAll(line, L"\\overline{b}=", L"");
    size_t i_p = line.find(L"i"), j_p = line.find(L"j"), k_p = line.find(L"k");
    if (i_p != wstring::npos && j_p != wstring::npos && k_p != wstring::npos) {
        auto trimP = [](wstring s) { 
            size_t i = 0; 
            while (i < s.length() && (s[i] == L' ' || s[i] == L'+')) i++; 
            return s.substr(i); };
        P = ParsePoly(line.substr(0, i_p)); 
        Q = ParsePoly(trimP(line.substr(i_p + 1, j_p - i_p - 1))); 
        R = ParsePoly(trimP(line.substr(j_p + 1, k_p - j_p - 1)));}
    else {
        wstringstream ss(line); wstring item;
        if (getline(ss, item, L';')) P = ParsePoly(item);
        if (getline(ss, item, L';')) Q = ParsePoly(item);
        if (getline(ss, item, L';')) R = ParsePoly(item);
    }
}
\end{listing}
\normalsize

\pagebreak%
\section{Метод FieldMathCalculations::ParsePolynomials}
\footnotesize
\begin{listing}{1}
/*************************************************************************
* Метод         : FieldMathCalculations::ParsePolynomials                *
*------------------------------------------------------------------------*
* Назначение    : Парсинг строковых выражений в структуру полиномов      *
* Входные       : a_str, b_str - строки формул полей A и B               *
* Выходные      : Нет (обновляет статические поля класса)                *
*------------------------------------------------------------------------*
* ОПИСАНИЕ ФУНКЦИОНИРОВАНИЯ:                                             *
* Метод преобразует текстовые представления векторных полей              *
* в математическую структуру Polynomial.                                   *
* *
* АЛГОРИТМ:                                                              *
* 1. Вызывает ExtractComponents для каждой строки поля (A и B).          *
* 2. ExtractComponents выполняет нормализацию и разделение на            *
* компоненты X, Y, Z.                                                    *
* 3. Результаты сохраняются в статические переменные класса.             *
*************************************************************************/
void FieldMathCalculations::ParsePolynomials(const wstring& a_str, 
                                             const wstring& b_str) {
    ExtractComponents(a_str, fieldA_X, fieldA_Y, fieldA_Z);
    ExtractComponents(b_str, fieldB_X, fieldB_Y, fieldB_Z);
}
\end{listing}
\normalsize

\pagebreak

\section{Метод FieldMathCalculations::EvaluateFieldA}
\footnotesize
\begin{listing}{1}
/*************************************************************************
* Метод         : FieldMathCalculations::EvaluateFieldA                  *
*------------------------------------------------------------------------*
* Назначение    : Вычисление вектора поля A в заданной точке             *
* Входные       : x, y, z - координаты точки                             *
* Выходные      : Vector3D - результирующий вектор поля в точке          *
*------------------------------------------------------------------------*
* ОПИСАНИЕ ФУНКЦИОНИРОВАНИЯ:                                             *
* Выполняет расчет численных значений компонент векторного поля A        *
* в указанной точке.                                                     *
* *
* АЛГОРИТМ:                                                              *
* 1. Обращается к компонентам fieldA_X, fieldA_Y, fieldA_Z.              *
* 2. Вызывает метод Evaluate для каждого компонента.                     *
* 3. Объединяет результаты в структуру Vector3D.                         *
*************************************************************************/
Vector3D FieldMathCalculations::EvaluateFieldA(double x, double y, double z) {
    return { fieldA_X.Evaluate(x, y, z), 
             fieldA_Y.Evaluate(x, y, z), 
             fieldA_Z.Evaluate(x, y, z) };
}
\end{listing}
\normalsize

\pagebreak

\section{Метод FieldMathCalculations::EvaluateFieldB}
\footnotesize
\begin{listing}{1}
/*************************************************************************
* Метод         : FieldMathCalculations::EvaluateFieldB                  *
*------------------------------------------------------------------------*
* Назначение    : Вычисление вектора поля B в заданной точке             *
* Входные       : x, y, z - координаты точки                             *
* Выходные      : Vector3D - результирующий вектор поля в точке          *
*------------------------------------------------------------------------*
* ОПИСАНИЕ ФУНКЦИОНИРОВАНИЯ:                                             *
* Аналогичен EvaluateFieldA, но работает с полем B, используя методы-    *
* геттеры для доступа к компонентам.                                     *
* *
* АЛГОРИТМ:                                                              *
* 1. Вызывает методы GetFieldBX, GetFieldBY, GetFieldBZ.                 *
* 2. Применяет метод Evaluate для каждого компонента.                    *
* 3. Возвращает вектор в виде структуры Vector3D.                        *
*************************************************************************/
Vector3D FieldMathCalculations::EvaluateFieldB(double x, double y, double z) {
    return { GetFieldBX().Evaluate(x, y, z), 
             GetFieldBY().Evaluate(x, y, z), 
             GetFieldBZ().Evaluate(x, y, z) };
}
\end{listing}
\normalsize

\pagebreak

\section{Метод FieldMathCalculations::Calc\_PartialDerivative}
\footnotesize
\begin{listing}{1}
/*************************************************************************
* Метод         : FieldMathCalculations::Calc_PartialDerivative          *
*------------------------------------------------------------------------*
* Назначение    : Численное вычисление частной производной функции       *
* Входные       : f - функция, variable - индекс переменной, x,y,z       *
* Выходные      : double - приближенное значение производной             *
*------------------------------------------------------------------------*
* ОПИСАНИЕ ФУНКЦИОНИРОВАНИЯ:                                             *
* Вычисляет производную функции f по переменной x, y или z               *
* методом центральной разности.                                          *
* *
* АЛГОРИТМ:                                                              *
* 1. Шаг дифференцирования h = 1e-4.                                     *
* 2. Вычисляет отношение: (f(x+h) - f(x-h)) / (2*h).                     *
*************************************************************************/
double FieldMathCalculations::Calc_PartialDerivative(FUNC3D f, int variable, 
                                                     double x, double y, 
                                                     double z) {
    double h = 1e-4;
    if (variable == 1) return (f(x + h, y, z) - f(x - h, y, z)) / (2 * h);
    if (variable == 2) return (f(x, y + h, z) - f(x, y - h, z)) / (2 * h);
    return (f(x, y, z + h) - f(x, y, z - h)) / (2 * h);
}
\end{listing}
\normalsize

\pagebreak

\section{Метод FieldMathCalculations::IntegrateQuarterDisk}
\footnotesize
\begin{listing}{1}
/*************************************************************************
* Метод         : FieldMathCalculations::IntegrateQuarterDisk            *
*------------------------------------------------------------------------*
* Назначение    : Численное интегрирование функции по четверти диска     *
* Входные       : func - подынтегральная функция, radius, steps          *
* Выходные      : double - значение интеграла                            *
*------------------------------------------------------------------------*
* ОПИСАНИЕ ФУНКЦИОНИРОВАНИЯ:                                             *
* Численное интегрирование двумерной функции по области четверти         *
* диска в полярных координатах (r, phi).                                 *
* *
* АЛГОРИТМ:                                                              *
* 1. Переход к полярным координатам: x = r*cos(phi), y = r*sin(phi). "phi"-   *
*обозначение полярного угла в системе полярных координат. *
* 2. Вычисление двойной суммы по сеткам r и phi с учетом якобиана (r).   *
*************************************************************************/
double FieldMathCalculations::IntegrateQuarterDisk(
    function<double(double, double)> func, double radius, int steps) {
    double dr = radius / steps, dphi = (M_PI / 2.0) / steps, sum = 0.0;
    for (int i = 0; i < steps; i++) {
        double r = (i + 0.5) * dr;
        for (int j = 0; j < steps; j++) {
            sum += func(r * cos((j + 0.5) * dphi), 
                        r * sin((j + 0.5) * dphi)) * r;
        }
    }
    return sum * dr * dphi;
}
\end{listing}
\normalsize

\pagebreak%
\section{Метод FieldMathCalculations::Calc\_Divergence}
\footnotesize
\begin{listing}{1}
/*************************************************************************
* Метод         : FieldMathCalculations::Calc_Divergence                 *
*------------------------------------------------------------------------*
* Назначение    : Численное вычисление дивергенции векторного поля       *
* Входные       : field - функция поля, x, y, z - координаты точки       *
* Выходные      : double - значение дивергенции в точке                  *
*------------------------------------------------------------------------*
* ОПИСАНИЕ ФУНКЦИОНИРОВАНИЯ:                                             *
* Вычисляет дивергенцию (div F) векторного поля, используя определение   *
* через частные производные компонент: div F = dP/dx + dQ/dy + dR/dz.     *
* *
* АЛГОРИТМ:                                                              *
* 1. Векторное поле раскладывается на скалярные функции P, Q, R.         *
* 2. Для каждой компоненты вычисляется соответствующая частная           *
* производная по своей переменной через метод Calc_PartialDerivative.    *
* 3. Результаты суммируются для получения итогового скалярного значения. *
*************************************************************************/
double FieldMathCalculations::Calc_Divergence(FUNC3D_VEC field, double x, 
                                              double y, double z) {
    FUNC3D P = [&](double px, double py, double pz) { 
        return field(px, py, pz).x; };
    FUNC3D Q = [&](double px, double py, double pz) { 
        return field(px, py, pz).y; };
    FUNC3D R = [&](double px, double py, double pz) { 
        return field(px, py, pz).z; };
    return Calc_PartialDerivative(P, 1, x, y, z) + 
           Calc_PartialDerivative(Q, 2, x, y, z) + 
           Calc_PartialDerivative(R, 3, x, y, z);
}
\end{listing}
\normalsize

\pagebreak

\section{Метод FieldMathCalculations::Calc\_Rotor}
\footnotesize
\begin{listing}{1}
/*************************************************************************
* Метод         : FieldMathCalculations::Calc_Rotor                      *
*------------------------------------------------------------------------*
* Назначение    : Численное вычисление ротора (curl) векторного поля     *
* Входные       : field - функция поля, x, y, z - координаты точки       *
* Выходные      : Vector3D - вектор ротора в точке                       *
*------------------------------------------------------------------------*
* ОПИСАНИЕ ФУНКЦИОНИРОВАНИЯ:                                             *
* Вычисляет вектор ротора (rot F) по формуле векторного произведения     *
* оператора набла на вектор поля.                                        *
* *
* АЛГОРИТМ:                                                              *
* 1. Метод вычисляет шесть частных производных компонент поля.           *
* 2. Формирует компоненты вектора ротора:                                *
* x: dR/dy - dQ/dz; y: dP/dz - dR/dx; z: dQ/dx - dP/dy.               *
*************************************************************************/
Vector3D FieldMathCalculations::Calc_Rotor(FUNC3D_VEC field, double x, 
                                           double y, double z) {
    FUNC3D P = [&](double px, double py, double pz) { 
        return field(px, py, pz).x; };
    FUNC3D Q = [&](double px, double py, double pz) { 
        return field(px, py, pz).y; };
    FUNC3D R = [&](double px, double py, double pz) { 
        return field(px, py, pz).z; };
    return {
        Calc_PartialDerivative(R, 2, x, y, z) - Calc_PartialDerivative(Q, 3, x, y, z),
        Calc_PartialDerivative(P, 3, x, y, z) - Calc_PartialDerivative(R, 1, x, y, z),
        Calc_PartialDerivative(Q, 1, x, y, z) - Calc_PartialDerivative(P, 2, x, y, z)
    };
}
\end{listing}
\normalsize

\pagebreak

\section{Метод FieldMathCalculations::DetermineFieldClass}
\footnotesize
\begin{listing}{1}
/*************************************************************************
* Метод         : FieldMathCalculations::DetermineFieldClass             *
*------------------------------------------------------------------------*
* Назначение    : Классификация типа векторного поля                     *
* Входные       : field - функция поля                                   *
* Выходные      : wstring - название класса поля                         *
*------------------------------------------------------------------------*
* ОПИСАНИЕ ФУНКЦИОНИРОВАНИЯ:                                             *
* Определяет тип поля (потенциальное, соленоидальное, гармоническое или  *
* общего вида) путем численной проверки свойств дивергенции и ротора.    *
*************************************************************************/
wstring FieldMathCalculations::DetermineFieldClass(FUNC3D_VEC field) {
    bool is_div = true, is_rot = true;
    for (double x = 0.2; x <= 0.8; x += 0.3) {
        for (double y = 0.2; y <= 0.8; y += 0.3) {
            if (abs(Calc_Divergence(field, x, y, 0.5)) > 1e-3) is_div = false;
            Vector3D r = Calc_Rotor(field, x, y, 0.5);
            if (abs(r.x) > 1e-3 || abs(r.y) > 1e-3 || abs(r.z) > 1e-3) 
                is_rot = false;
        }
    }
    if (is_div && is_rot) return L"Гармоническое";
    return (is_rot ? L"Потенциальное" : 
           (is_div ? L"Соленоидальное" : L"Общего вида"));
}
\end{listing}
\normalsize

\pagebreak

\section{Метод FieldMathCalculations::Calc\_Flux\_GaussOstrogradsky}
\footnotesize
\begin{listing}{1}
/*************************************************************************
* Метод         : FieldMathCalculations::Calc_Flux_GaussOstrogradsky     *
*------------------------------------------------------------------------*
* Назначение    : Расчет потока через теорему Остроградского-Гаусса      *
* Входные       : Нет                                                    *
* Выходные      : double - значение потока                               *
*------------------------------------------------------------------------*
* ОПИСАНИЕ ФУНКЦИОНИРОВАНИЯ:                                             *
* Реализует вычисление потока векторного поля через замкнутую поверхность *
* путем сведения поверхностного интеграла к объемному по теореме Гаусса. *
*************************************************************************/
double FieldMathCalculations::Calc_Flux_GaussOstrogradsky() {
    double Pc = 0.0; int st = 25;
    double d_r = 1.0 / st, d_p = (M_PI / 2.0) / st, d_t = (M_PI / 2.0) / st;
    for (int i = 0; i < st; ++i) {
        for (int j = 0; j < st; ++j) {
            for (int k = 0; k < st; ++k) {
                double r = (i + 0.5) * d_r, p = (j + 0.5) * d_p, t = (k + 0.5) * d_t;
                double x = r * cos(t) * cos(p), y = r * cos(t) * sin(p), z = r * sin(t);
                Pc += Calc_Divergence(EvaluateFieldA, x, y, z) 
	      	* (r * r * cos(t)) * d_r * d_p * d_t;
            }
        }
    }
    double P_z0 = -IntegrateQuarterDisk([](double x, double y) { 
        return EvaluateFieldA(x, y, 0.0).z; });
    double P_x0 = -IntegrateQuarterDisk([](double y, double z) { 
        return EvaluateFieldA(0.0, y, z).x; });
    double P_y0 = -IntegrateQuarterDisk([](double x, double z) { 
        return EvaluateFieldA(x, 0.0, z).y; });
    return Pc - (P_z0 + P_x0 + P_y0);
}
\end{listing}
\normalsize

\pagebreak

\section{Метод FieldMathCalculations::Calc\_Flux\_Direct}
\footnotesize
\begin{listing}{1}
/*************************************************************************
* Метод         : FieldMathCalculations::Calc_Flux_Direct                *
*------------------------------------------------------------------------*
* Назначение    : Прямое вычисление потока векторного поля              *
* Входные       : Нет                                                    *
* Выходные      : double - значение потока                               *
*------------------------------------------------------------------------*
* ОПИСАНИЕ ФУНКЦИОНИРОВАНИЯ:                                             *
* Выполняет прямое интегрирование векторного поля по поверхности сферы   *
* в первой октанте (четверть диска).                                     *
*************************************************************************/
double FieldMathCalculations::Calc_Flux_Direct() {
    return IntegrateQuarterDisk([](double x, double y) {
        double z = sqrt(fmax(0.0, 1.0 - x * x - y * y));
        if (z < 1e-5) return 0.0;
        Vector3D a = EvaluateFieldA(x, y, z);
        return (a.x * x + a.y * y + a.z * z) / z;
    });
}
\end{listing}
\normalsize

\pagebreak
%
\section{Метод FieldMathCalculations::Calc\_Circulation\_Stokes}
\footnotesize
\begin{listing}{1}
/*************************************************************************
* Метод         : FieldMathCalculations::Calc_Circulation_Stokes         *
*------------------------------------------------------------------------*
* Назначение    : Расчет циркуляции векторного поля через теорему Стокса *
* Входные       : Нет                                                    *
* Выходные      : double - значение циркуляции                           *
*------------------------------------------------------------------------*
* ОПИСАНИЕ ФУНКЦИОНИРОВАНИЯ:                                             *
* Вычисляет циркуляцию поля по контуру, используя теорему Стокса (переход*
* от криволинейного интеграла к интегралу от ротора по поверхности).     *
* *
* АЛГОРИТМ:                                                              *
* 1. Интегрирование ротора поля по поверхности через IntegrateQuarterDisk. *
* 2. Вычисление проекции вектора ротора на нормаль к поверхности.         *
*************************************************************************/
double FieldMathCalculations::Calc_Circulation_Stokes() {
    return IntegrateQuarterDisk([](double x, double y) {
        double z = sqrt(fmax(0.0, 1.0 - x * x - y * y));
        if (z < 1e-5) return 0.0;
        Vector3D rot = Calc_Rotor(EvaluateFieldA, x, y, z);
        return (rot.x * x + rot.y * y + rot.z * z) / z;
    });
}
\end{listing}
\normalsize

\pagebreak

\section{Метод FieldMathCalculations::Calc\_Circulation\_Direct}
\footnotesize
\begin{listing}{1}
/*************************************************************************
* Метод         : FieldMathCalculations::Calc_Circulation_Direct         *
*------------------------------------------------------------------------*
* Назначение    : Прямое вычисление циркуляции                          *
* Входные       : Нет                                                    *
* Выходные      : double - значение циркуляции                           *
*------------------------------------------------------------------------*
* ОПИСАНИЕ ФУНКЦИОНИРОВАНИЯ:                                             *
* Выполняет непосредственное интегрирование поля вдоль граничных         *
* отрезков в первой октанте (три стороны контура).                       *
* *
* АЛГОРИТМ:                                                              *
* 1. Параметризация трех сторон контура.                                 *
* 2. Численное суммирование скалярного произведения поля на касательный  *
* вектор (dl) с шагом dphi.                                              *
*************************************************************************/
double FieldMathCalculations::Calc_Circulation_Direct() {
    int steps = 100; 
    double dphi = (M_PI / 2.0) / steps, sum = 0.0;
    for (int i = 0; i < steps; i++) {
        double phi = (i + 0.5) * dphi;
        Vector3D a1 = EvaluateFieldA(cos(phi), sin(phi), 0.0);
        sum += (a1.x * (-sin(phi)) + a1.y * cos(phi));
        Vector3D a2 = EvaluateFieldA(0.0, cos(phi), sin(phi));
        sum += (a2.y * (-sin(phi)) + a2.z * cos(phi));
        Vector3D a3 = EvaluateFieldA(sin(phi), 0.0, cos(phi));
        sum += (a3.z * (-sin(phi)) + a3.x * cos(phi));
    }
    return sum * dphi;
}
\end{listing}
\normalsize

\pagebreak

\section{Метод FieldMathCalculations::GetAnalyticalPotential}
\footnotesize
\begin{listing}{1}
/*************************************************************************
* Метод         : FieldMathCalculations::GetAnalyticalPotential          *
*------------------------------------------------------------------------*
* Назначение    : Аналитическое нахождение потенциала поля B             *
* Входные       : Нет                                                    *
* Выходные      : Polynomial - результирующая функция потенциала         *
*------------------------------------------------------------------------*
* ОПИСАНИЕ ФУНКЦИОНИРОВАНИЯ:                                             *
* Вычисляет скалярный потенциал U(x,y,z), такой что grad(U) = B.         *
* *
* АЛГОРИТМ:                                                              *
* 1. Интегрирование компонент (Bx по x, By по y, Bz по z).               *
* 2. Использование метода SetVariableToZero для исключения перекрестных  *
* членов, чтобы избежать дублирования данных при восстановлении.        *
* 3. Алгебраическое сложение полученных функций (Add).                   *
*************************************************************************/
Polynomial FieldMathCalculations::GetAnalyticalPotential() {
    Polynomial U1 = GetFieldBX().SetVariableToZero(false, true, true)
                                .IntegrateSymbolically(0);
    Polynomial U2 = GetFieldBY().SetVariableToZero(false, false, true)
                                .IntegrateSymbolically(1);
    Polynomial U3 = GetFieldBZ().SetVariableToZero(false, false, false)
                                .IntegrateSymbolically(2);
    Polynomial U; U.Add(U1); U.Add(U2); U.Add(U3);
    return U;
}
\end{listing}
\normalsize

\pagebreak

\section{Метод FieldMathCalculations::VerifyPotentialCorrectness}
\footnotesize
\begin{listing}{1}
/*************************************************************************
* Метод         : FieldMathCalculations::VerifyPotentialCorrectness      *
*------------------------------------------------------------------------*
* Назначение    : Верификация точности найденного потенциала             *
* Входные       : Нет                                                    *
* Выходные      : bool - статус корректности                             *
*------------------------------------------------------------------------*
* ОПИСАНИЕ ФУНКЦИОНИРОВАНИЯ:                                             *
* Проверяет условие grad(U) = B в контрольных точках сетки.              *
* *
* АЛГОРИТМ:                                                              *
* 1. В контрольных точках вычисляются частные производные от U.          *
* 2. Сравнение градиента с вектором поля B. Если разница превышает       *
* порог (0.05) в любой точке, возвращается false.                       *
*************************************************************************/
bool FieldMathCalculations::VerifyPotentialCorrectness() {
    Polynomial U = GetAnalyticalPotential();
    FUNC3D U_eval = [&](double px, double py, double pz) { 
        return U.Evaluate(px, py, pz); 
    };
    for (double x = 0.2; x <= 0.8; x += 0.3) {
        for (double y = 0.2; y <= 0.8; y += 0.3) {
            double du_dx = Calc_PartialDerivative(U_eval, 1, x, y, 0.5);
            double du_dy = Calc_PartialDerivative(U_eval, 2, x, y, 0.5);
            Vector3D valB = EvaluateFieldB(x, y, 0.5);
            if (abs(du_dx - valB.x) > 0.05 || abs(du_dy - valB.y) > 0.05) 
                return false;
        }
    }
    return true;
}
\end{listing}
\normalsize

\pagebreak

\section{Метод FieldMathCalculations::CheckRotorBZero}
\footnotesize
\begin{listing}{1}
/*************************************************************************
* Метод         : FieldMathCalculations::CheckRotorBZero                 *
*------------------------------------------------------------------------*
* Назначение    : Проверка потенциальности поля B (rot B = 0)            *
* Входные       : Нет                                                    *
* Выходные      : bool - результат проверки                              *
*------------------------------------------------------------------------*
* ОПИСАНИЕ ФУНКЦИОНИРОВАНИЯ:                                             *
* Численная проверка условия отсутствия вихрей в поле B.                 *
* *
* АЛГОРИТМ:                                                              *
* 1. Сканирование области по сетке (x, y).                               *
* 2. Вычисление ротора в каждой точке.                                   *
* 3. Если модуль ротора превышает 1e-3, поле считается непотенциальным.  *
*************************************************************************/
bool FieldMathCalculations::CheckRotorBZero() {
    for (double x = 0.2; x <= 0.8; x += 0.3) {
        for (double y = 0.2; y <= 0.8; y += 0.3) {
            Vector3D r = Calc_Rotor(EvaluateFieldB, x, y, 0.5);
            if (abs(r.x) > 1e-3 || abs(r.y) > 1e-3 || abs(r.z) > 1e-3) 
                return false;
        }
    }
    return true;
}
\end{listing}
\normalsize

\pagebreak
%
\section{Функция GetLoadFileNameExplorer}
\footnotesize
\begin{listing}{1}
/*************************************************************************
* Функция       : GetLoadFileNameExplorer                                *
*------------------------------------------------------------------------*
* Назначение    : Диалог выбора файла для открытия                       *
* Входные       : hWnd - дескриптор окна, outPath - путь к файлу         *
* Выходные      : bool - успех операции                                  *
*------------------------------------------------------------------------*
* ОПИСАНИЕ ФУНКЦИОНИРОВАНИЯ:                                             *
* Обертка над Win32 API для вызова стандартного диалога открытия файла.  *
*************************************************************************/
bool GetLoadFileNameExplorer(HWND hWnd, wstring& outPath) {
    wchar_t filename[MAX_PATH] = L""; 
    OPENFILENAMEW ofn; 
    ZeroMemory(&ofn, sizeof(ofn));
    ofn.lStructSize = sizeof(ofn); 
    ofn.hwndOwner = hWnd;
    ofn.lpstrFilter = L"Текстовые файлы (*.txt)\0*.txt\0Все файлы (*.*)\0*.*\0";
    ofn.lpstrFile = filename; 
    ofn.nMaxFile = MAX_PATH;
    ofn.Flags = OFN_FILEMUSTEXIST | OFN_PATHMUSTEXIST;
    if (GetOpenFileNameW(&ofn)) { 
        outPath = filename; 
        return true; 
    }
    return false;
}
\end{listing}
\normalsize

\pagebreak

\section{Функция GetSaveFileNameExplorer}
\footnotesize
\begin{listing}{1}
/*************************************************************************
* Функция       : GetSaveFileNameExplorer                                *
*------------------------------------------------------------------------*
* Назначение    : Диалог выбора файла для сохранения                     *
* Входные       : hWnd - дескриптор окна, outPath - путь к файлу         *
* Выходные      : bool - успех операции                                  *
*------------------------------------------------------------------------*
* ОПИСАНИЕ ФУНКЦИОНИРОВАНИЯ:                                             *
* Обертка над Win32 API для вызова стандартного диалога сохранения файла. *
*************************************************************************/
bool GetSaveFileNameExplorer(HWND hWnd, wstring& outPath) {
    wchar_t filename[MAX_PATH] = L"output.txt"; 
    OPENFILENAMEW ofn; 
    ZeroMemory(&ofn, sizeof(ofn));
    ofn.lStructSize = sizeof(ofn); 
    ofn.hwndOwner = hWnd;
    ofn.lpstrFilter = L"Текстовые файлы (*.txt)\0*.txt\0";
    ofn.lpstrFile = filename; 
    ofn.nMaxFile = MAX_PATH;
    ofn.Flags = OFN_OVERWRITEPROMPT | OFN_PATHMUSTEXIST;
    if (GetSaveFileNameW(&ofn)) { 
        outPath = filename; 
        return true; 
    }
    return false;
}
\end{listing}
\normalsize

\pagebreak

\section{Функция ImportDataFromFile}
\footnotesize
\begin{listing}{1}
/*************************************************************************
* Функция       : ImportDataFromFile                                     *
*------------------------------------------------------------------------*
* Назначение    : Импорт и парсинг данных из текстового файла            *
* Входные       : filename - путь к файлу                                *
* Выходные      : bool - успех импорта                                   *
*------------------------------------------------------------------------*
* ОПИСАНИЕ ФУНКЦИОНИРОВАНИЯ:                                             *
* Читает вариант задания и формулы полей, выполняет их первичный парсинг. *
*************************************************************************/
bool ImportDataFromFile(const wstring& filename) {
    ifstream file(filename);
    if (!file.is_open()) return false;
    string line;
    if (getline(file, line)) { 
        try { g_VariantNum = stoi(line); } 
        catch (...) { return false; } 
    }
    if (getline(file, line)) 
        g_FieldAStr = wstring(line.begin(), line.end());
    if (getline(file, line)) 
        g_FieldBStr = wstring(line.begin(), line.end());
    file.close();
    if (g_VariantNum <= 0 || g_FieldAStr.empty() || g_FieldBStr.empty()) 
        return false;
    FieldMathCalculations::ParsePolynomials(g_FieldAStr, g_FieldBStr);
    return true;
}
\end{listing}
\normalsize

\pagebreak

\section{Функция ExportDataToFile}
\footnotesize
\begin{listing}{1}
/*************************************************************************
* Функция       : ExportDataToFile                                       *
*------------------------------------------------------------------------*
* Назначение    : Экспорт отчета в файл (UTF-8 с BOM)                    *
* Входные       : filename - путь к файлу                                *
* Выходные      : bool - успех записи                                    *
*------------------------------------------------------------------------*
* ОПИСАНИЕ ФУНКЦИОНИРОВАНИЯ:                                             *
* Конвертирует отчет в UTF-8 и записывает в файл с добавлением BOM-метки. *
*************************************************************************/
bool ExportDataToFile(const wstring& filename) {
    wstring report = BuildReportText();
    if (report.empty() || g_VariantNum == 0) return false;
    int s_n = WideCharToMultiByte(CP_UTF8, 0, &report[0], 
                                  (int)report.size(), NULL, 0, NULL, NULL);
    string utf8_str(s_n, 0);
    WideCharToMultiByte(CP_UTF8, 0, &report[0], (int)report.size(), 
                        &utf8_str[0], s_n, NULL, NULL);                  
    ofstream file(filename, ios::binary);
    if (!file.is_open()) return false;
    const unsigned char bom[] = { 0xEF, 0xBB, 0xBF };
    file.write((const char*)bom, sizeof(bom)); 
    file.write(utf8_str.c_str(), utf8_str.size()); 
    file.close();
    return true;
}
\end{listing}
\normalsize

\pagebreak

\section{Функция BuildInputDataText}
\footnotesize
\begin{listing}{1}
/*************************************************************************
* Функция       : BuildInputDataText                                     *
*------------------------------------------------------------------------*
* Назначение    : Формирование текстового блока исходных данных          *
* Входные       : Нет                                                    *
* Выходные      : wstring - отформатированный текст                      *
*------------------------------------------------------------------------*
* ОПИСАНИЕ ФУНКЦИОНИРОВАНИЯ:                                             *
* Создает текстовое представление входных данных для отображения в GUI.  *
*************************************************************************/
wstring BuildInputDataText() {
    wstringstream ss;
    ss << L"====================================================================\n";
    ss << L"        ИСХОДНЫЕ ДАННЫЕ ВАРИАНТА\n";
    ss << L"====================================================================\n\n";
    ss << L"  ► Загружен вариант №:  " << g_VariantNum << L"\n\n";
    ss << L"  ► Векторное поле ā (исследуется на поток и циркуляцию):\n";
    ss << L"    ā = " << g_FieldAStr << L"\n\n";
    ss << L"  ► Векторное поле b̄ (исследуется на потенциальность):\n";
    ss << L"    b̄ = " << g_FieldBStr << L"\n\n";
    ss << L"--------------------------------------------------------------------\n";
    ss << L"  Всё готово! Нажмите кнопку «Расчёт данных» для получения отчета.";
    return ss.str();
}
\end{listing}
\normalsize

\pagebreak%
\section{Функция BuildReportText}
\footnotesize
\begin{listing}{1}
/*************************************************************************
* Функция       : BuildReportText                                        *
*------------------------------------------------------------------------*
* Назначение    : Формирование полного отчета по расчетам                *
* Входные       : Нет                                                    *
* Выходные      : wstring - текст итогового отчета                       *
*------------------------------------------------------------------------*
* ОПИСАНИЕ ФУНКЦИОНИРОВАНИЯ:                                             *
* Собирает результаты вычислений (потоки, циркуляции, потенциал)          *
* в единую строку для вывода в интерфейс или сохранения в файл.          *
*************************************************************************/
wstring BuildReportText() { wstringstream rep; rep << fixed << setprecision(4); 
    Polynomial U = FieldMathCalculations::GetAnalyticalPotential();
    rep << L"========== РАССЧИТАННЫЕ ДАННЫЕ (ВАР. №" << g_VariantNum << L") ======\n\n"
        << L"► 1. КЛАССИФИКАЦИЯ ПОЛЯ ā:\n   • Тип поля: " 
        << FieldMathCalculations::DetermineFieldClass
	( FieldMathCalculations::EvaluateFieldA) << L"\n\n"
        << L"► 2. ВЫЧИСЛЕНИЕ ПОТОКА ПОВЕРХНОСТИ S (ПОЛЕ ā):\n"
        << L"   • 2.1 Метод Остроградского-Гаусса: P = " 
        << FieldMathCalculations::Calc_Flux_GaussOstrogradsky() << L"\n"
        << L"   • 2.2 Метод прямого интегрирования: P = " 
        << FieldMathCalculations::Calc_Flux_Direct() << L"\n\n"
        << L"► 3. ВЫЧИСЛЕНИЕ ЦИРКУЛЯЦИИ ПО КОНТУРУ S (ПОЛЕ ā):\n"
        << L"   • 3.1 Метод теоремы Стокса: C = " 
        << FieldMathCalculations::Calc_Circulation_Stokes() << L"\n"
        << L"   • 3.2 Метод криволинейного интеграл: C = " 
        << FieldMathCalculations::Calc_Circulation_Direct() << L"\n\n"
        << L"► 4. ИССЛЕДОВАНИЕ ПОЛЯ b̄:\n" << L"   • 4.1 Проверка потенциальности: " 
        << (FieldMathCalculations::CheckRotorBZero() ? 
            L"поле потенциально" : L"поле не потенциально") << L"\n"
        << L"   • 4.2 Потенциал U(x, y, z) = " << U.ToWString() << L"\n"
        << L"U(1, 1, 1) = " << U.Evaluate(1, 1, 1) << L"\n"<< L" • 4.3 Проверка градиента: " 
        << (FieldMathCalculations::VerifyPotentialCorrectness() ? 
            L"успешна" : L"не успешна") << L"\n" << L"   • 4.4 Работа сил поля W = " 
        << U.Evaluate(1, 1, 1) - U.Evaluate(0, 0, 0) << L"\n\n"
        << L"========================================================="; return rep.str(); }
\end{listing}
\normalsize

\pagebreak

\section{Процедура CreateInterfaceButtons}
\footnotesize
\begin{listing}{1}
/*************************************************************************
* Процедура     : CreateInterfaceButtons                                 *
*------------------------------------------------------------------------*
* Назначение    : Создание кнопок управления интерфейсом                 *
* Входные       : hWnd - окно, hInstance - экземпляр приложения          *
* Выходные      : Нет                                                    *
*------------------------------------------------------------------------*
* ОПИСАНИЕ ФУНКЦИОНИРОВАНИЯ:                                             *
* Инициализирует элементы управления (кнопки) в рабочей области окна.    *
*************************************************************************/
void CreateInterfaceButtons(HWND hWnd, HINSTANCE hInstance) {
    int y1 = 45, w = 185, h = 35, step = 190;
    CreateWindowW(L"BUTTON", L"1. Импорт данных", WS_VISIBLE | WS_CHILD | BS_PUSHBUTTON, 
                  20 + 0 * step, y1, w, h, hWnd, (HMENU)ID_BTN_IMPORT, 
                  hInstance, NULL);
    CreateWindowW(L"BUTTON", L"2. Показать данные", WS_VISIBLE | WS_CHILD | BS_PUSHBUTTON, 
                  20 + 1 * step, y1, w, h, hWnd, (HMENU)ID_BTN_SHOW, 
                  hInstance, NULL);
    CreateWindowW(L"BUTTON", L"3. Расчёт данных", WS_VISIBLE | WS_CHILD | BS_PUSHBUTTON, 
                  20 + 2 * step, y1, w, h, hWnd, (HMENU)ID_BTN_CALC, 
                  hInstance, NULL);
    CreateWindowW(L"BUTTON", L"4. Визуализация Поля А", WS_VISIBLE | WS_CHILD | BS_PUSHBUTTON, 
                  20 + 3 * step, y1, w, h, hWnd, (HMENU)ID_BTN_VIS_A, 
                  hInstance, NULL);
    CreateWindowW(L"BUTTON", L"5. Визуализация Поля Б", WS_VISIBLE | WS_CHILD | BS_PUSHBUTTON, 
                  20 + 4 * step, y1, w, h, hWnd, (HMENU)ID_BTN_VIS_B, hInstance, NULL);
    CreateWindowW(L"BUTTON", L"6. Экспорт в TXT", WS_VISIBLE | WS_CHILD | BS_PUSHBUTTON, 
                  20 + 5 * step, y1, w, h, hWnd, (HMENU)ID_BTN_EXPORT, hInstance, NULL);
}
\end{listing}
\normalsize

\pagebreak

\section{Процедура DrawBackground}
\footnotesize
\begin{listing}{1}
/*************************************************************************
* Процедура     : DrawBackground                                         *
*------------------------------------------------------------------------*
* Назначение    : Отрисовка фона и информации о разработчиках            *
* Входные       : hdc - контекст устройства, rect - область отрисовки    *
* Выходные      : Нет                                                    *
*------------------------------------------------------------------------*
* ОПИСАНИЕ ФУНКЦИОНИРОВАНИЯ:                                             *
* Заливает фон цветом и выводит строку с информацией о разработчиках.    *
*************************************************************************/
void DrawBackground(HDC hdc, RECT rect) {
    HBRUSH hBgBrush = CreateSolidBrush(RGB(240, 244, 248)); 
    FillRect(hdc, &rect, hBgBrush); 
    DeleteObject(hBgBrush);
    HFONT hFontDev = CreateFontW(18, 0, 0, 0, FW_SEMIBOLD, 0, 0, 0, DEFAULT_CHARSET,
		 OUT_OUTLINE_PRECIS, CLIP_DEFAULT_PRECIS,
 		 CLEARTYPE_QUALITY, VARIABLE_PITCH, L"Segoe UI");
    HFONT hOld = (HFONT)SelectObject(hdc, hFontDev); 
    SetBkMode(hdc, TRANSPARENT); SetTextColor(hdc, RGB(0, 0, 0));
    wstring devInfo = L"Разработчики: Варфоломеев, Великоцкая, Нагороднюк "
                      L"  |   Группа: М3О-225БВ-24";
    RECT devRect = { 25, 10, rect.right - 20, 35 }; 
    DrawTextW(hdc, devInfo.c_str(), -1, &devRect, 
              DT_LEFT | DT_VCENTER | DT_SINGLELINE);     
    SelectObject(hdc, hOld); DeleteObject(hFontDev);
}
\end{listing}
\normalsize

\pagebreak

\section{Процедура DrawStatusAndCard}
\footnotesize
\begin{listing}{1}
/*************************************************************************
* Процедура     : DrawStatusAndCard                                      *
*------------------------------------------------------------------------*
* Назначение    : Отрисовка панели статуса и центральной "карточки"      *
* Входные       : hdc - контекст устройства, rect - область отрисовки    *
* Выходные      : Нет                                                    *
*------------------------------------------------------------------------*
* ОПИСАНИЕ ФУНКЦИОНИРОВАНИЯ:                                             *
* Рисует полосу статуса и прямоугольник (карточку) для вывода контента.   *
*************************************************************************/
void DrawStatusAndCard(HDC hdc, RECT rect) {
    RECT sRect = { 20, 100, rect.right - 20, 130 }; 
    HBRUSH hStatB = CreateSolidBrush(RGB(225, 235, 245)); 
    FillRect(hdc, &sRect, hStatB); DeleteObject(hStatB);
    HFONT hFontS = CreateFontW(16, 0, 0, 0, FW_BOLD, 0, 0, 0, DEFAULT_CHARSET,
		 OUT_OUTLINE_PRECIS, CLIP_DEFAULT_PRECIS,
		 CLEARTYPE_QUALITY, VARIABLE_PITCH, L"Segoe UI");
    HFONT hOld = (HFONT)SelectObject(hdc, hFontS); 
    SetTextColor(hdc, RGB(0, 0, 0));
    RECT sText = { 35, 105, rect.right - 35, 125 }; 
    DrawTextW(hdc, g_StatusMessage.c_str(), -1, &sText, 
              DT_LEFT | DT_VCENTER | DT_SINGLELINE);
    RECT cR = { 20, 145, rect.right - 20, rect.bottom - 20 }; 
    RECT shR = { cR.left + 3, cR.top + 3, cR.right + 3, cR.bottom + 3 };
    HBRUSH hShB = CreateSolidBrush(RGB(220, 225, 230)); 
    FillRect(hdc, &shR, hShB); DeleteObject(hShB);
    HBRUSH hCardB = CreateSolidBrush(RGB(255, 255, 255)); 
    HPEN hCardP = CreatePen(PS_SOLID, 1, RGB(200, 205, 215));
    HPEN hOldP = (HPEN)SelectObject(hdc, hCardP); 
    HBRUSH hOldB = (HBRUSH)SelectObject(hdc, hCardB);
    Rectangle(hdc, cR.left, cR.top, cR.right, cR.bottom);
    SelectObject(hdc, hOldB); SelectObject(hdc, hOldP); 
    DeleteObject(hCardB); DeleteObject(hCardP); 
    SelectObject(hdc, hOld); DeleteObject(hFontS);
}
\end{listing}
\normalsize

\pagebreak

\section{Процедура HandlePaint}
\footnotesize
\begin{listing}{1}
/*************************************************************************
* Процедура     : HandlePaint                                            *
*------------------------------------------------------------------------*
* Назначение    : Обработчик события перерисовки окна (WM_PAINT)         *
* Входные       : hWnd - дескриптор окна                                 *
* Выходные      : Нет                                                    *
*------------------------------------------------------------------------*
* ОПИСАНИЕ ФУНКЦИОНИРОВАНИЯ:                                             *
* Вызывает функции отрисовки интерфейса и выводит контент в зависимости    *
* от выбранного пользователем режима (вид 0-4).                          *
*************************************************************************/
void HandlePaint(HWND hWnd) {
    PAINTSTRUCT ps; HDC hdc = BeginPaint(hWnd, &ps); RECT r; GetClientRect(hWnd, &r);
    DrawBackground(hdc, r); DrawStatusAndCard(hdc, r);
    HFONT hF = CreateFontW(20, 0, 0, 0, FW_NORMAL, 0, 0, 0, DEFAULT_CHARSET,
        OUT_OUTLINE_PRECIS,  CLIP_DEFAULT_PRECIS, CLEARTYPE_QUALITY, 
        VARIABLE_PITCH, L"Segoe UI");
    HFONT hOld = (HFONT)SelectObject(hdc, hF); SetTextColor(hdc, RGB(0, 0, 0)); 
    RECT tR = { 40, 165, r.right - 40, r.bottom - 40 };
    if (g_ActiveView == 1) { wstring t = BuildInputDataText(); 
        DrawTextW(hdc, t.c_str(), -1, &tR, DT_LEFT | DT_WORDBREAK);}
    else if (g_ActiveView == 2) { wstring t = BuildReportText(); 
        DrawTextW(hdc, t.c_str(), -1, &tR, DT_LEFT | DT_WORDBREAK);}
    else if (g_ActiveView == 3 || g_ActiveView == 4) DrawVisualization(hdc, tR);
    else if (g_ActiveView == 0) { wstring t = L"\n\n ПРОГРАММА РАСЧЕТА ПАРАМЕТРОВ "
                    L"ВЕКТОРНЫХ ПОЛЕЙ\n       -------------------------"
                    L"-------------------\n\n"
                    L"       Для работы с программой:\n\n"
                    L"       1. Нажмите кнопку «1. Импорт данных» сверху\n"
                    L"       2. Выберите .txt файл с вариантом задания\n"
                    L"       3. Нажмите «3. Расчёт данных» для отчета\n"
                    L"       4. Используйте кнопки визуализации.\n"
                    L"       5. Сохраните результат (кнопка 6).\n";
        DrawTextW(hdc, t.c_str(), -1, &tR, DT_LEFT | DT_WORDBREAK);}
    SelectObject(hdc, hOld); DeleteObject(hF); EndPaint(hWnd, &ps);
}
\end{listing}
\normalsize
\pagebreak%
\section{Процедура HandleCmdImportShow}
\footnotesize
\begin{listing}{1}
/*************************************************************************
* Процедура     : HandleCmdImportShow                                    *
*------------------------------------------------------------------------*
* Назначение    : Обработка команд импорта и отображения данных          *
* Входные       : hWnd - дескриптор окна, wmId - идентификатор команды   *
* Выходные      : Нет                                                    *
*------------------------------------------------------------------------*
* ОПИСАНИЕ ФУНКЦИОНИРОВАНИЯ:                                             *
* Обрабатывает нажатия кнопок импорта и показа данных, управляет         *
* состоянием интерфейса и сбросом предыдущих результатов.                *
*************************************************************************/
void HandleCmdImportShow(HWND hWnd, int wmId) {
    if (wmId == ID_BTN_IMPORT) { wstring f;
        if (GetLoadFileNameExplorer(hWnd, f)) {
            if (ImportDataFromFile(f)) {
                g_StatusMessage = L"✔ Успех! Вариант №" + to_wstring(g_VariantNum);
                g_ActiveView = 1; g_IsCalculated = false;
                if (g_hPlotBitmapA) { DeleteObject(g_hPlotBitmapA); g_hPlotBitmapA = NULL;}
                if (g_hPlotBitmapB) { DeleteObject(g_hPlotBitmapB); 
		g_hPlotBitmapB = NULL; }
	} else ShowError(hWnd, ErrParse);}
}
    else if (wmId == ID_BTN_SHOW) { if (g_VariantNum == 0) ShowError(hWnd, WarnNoData);
        else { g_StatusMessage = L"► Исходные данные ВАР №" + to_wstring(g_VariantNum); 
            g_ActiveView = 1; }
    }
}
\end{listing}
\normalsize

\pagebreak

\section{Процедура HandleCmdCalcExport}
\footnotesize
\begin{listing}{1}
/*************************************************************************
* Процедура     : HandleCmdCalcExport                                    *
*------------------------------------------------------------------------*
* Назначение    : Обработка команд расчета и экспорта                    *
* Входные       : hWnd - дескриптор окна, wmId - идентификатор команды   *
* Выходные      : Нет                                                    *
*------------------------------------------------------------------------*
* ОПИСАНИЕ ФУНКЦИОНИРОВАНИЯ:                                             *
* Выполняет расчет параметров поля и экспорт отчета в файл.              *
*************************************************************************/
void HandleCmdCalcExport(HWND hWnd, int wmId) {
    if (wmId == ID_BTN_CALC) {
        if (g_VariantNum == 0) ShowError(hWnd, WarnNoData);
        else { g_StatusMessage = L"✔ Расчёт выполнен!"; g_ActiveView = 2; 
            g_IsCalculated = true;}
    }
    else if (wmId == ID_BTN_EXPORT) {
        if (!g_IsCalculated) { ShowError(hWnd, WarnNoCalc); 
            return;}
        wstring s;
        if (GetSaveFileNameExplorer(hWnd, s)) {
            if (ExportDataToFile(s)) 
                g_StatusMessage = L"✔ Отчёт сохранён!";
            else 
                ShowError(hWnd, ErrWrite);}
    }
}
\end{listing}
\normalsize

\pagebreak

\section{Процедура HandleCmdVis}
\footnotesize
\begin{listing}{1}
/*************************************************************************
* Процедура     : HandleCmdVis                                           *
*------------------------------------------------------------------------*
* Назначение    : Обработка визуализации полей                           *
* Входные       : hWnd - дескриптор окна, wmId - идентификатор команды   *
* Выходные      : Нет                                                    *
*------------------------------------------------------------------------*
* ОПИСАНИЕ ФУНКЦИОНИРОВАНИЯ:                                             *
* Инициирует рендеринг графиков через Wolfram (с проверкой готовности),  *
* управляет состоянием отображения и выводит ошибки при сбоях.           *
*************************************************************************/
void HandleCmdVis(HWND hWnd, int wmId) {
    if (wmId == ID_BTN_VIS_A || wmId == ID_BTN_VIS_B) {
        if (!g_IsCalculated) { ShowError(hWnd, WarnNoCalc); return; }
        bool isA = (wmId == ID_BTN_VIS_A); 
        HBITMAP& bmp = isA ? g_hPlotBitmapA : g_hPlotBitmapB;
        if (!bmp) { g_StatusMessage = isA ? L"⏳ Рендеринг Поля А..." : 
                                    L"⏳ Рендеринг Поля Б..."; g_ActiveView = 5; 
            InvalidateRect(hWnd, NULL, TRUE); UpdateWindow(hWnd); ErrCode err;
            if (!GenerateWolframPlot(isA, err)) {
                g_StatusMessage = L"❌ Ошибка рендеринга!"; g_ActiveView = 2;
                InvalidateRect(hWnd, NULL, TRUE); UpdateWindow(hWnd); ShowError(hWnd, err); 
                return;}
        }
        g_StatusMessage = isA ? L"✔ Визуализация Поля А готова!" : 
                                L"✔ Визуализация Поля Б готова!";
        g_ActiveView = isA ? 3 : 4;}
}
\end{listing}
\normalsize

\pagebreak

\section{Процедура HandleWmCommand}
\footnotesize
\begin{listing}{1}
/*************************************************************************
* Процедура     : HandleWmCommand                                        *
*------------------------------------------------------------------------*
* Назначение    : Диспетчеризация команд меню и кнопок                   *
* Входные       : hWnd - окно, wParam - параметры сообщения              *
* Выходные      : Нет                                                    *
*------------------------------------------------------------------------*
* ОПИСАНИЕ ФУНКЦИОНИРОВАНИЯ:                                             *
* Центральный обработчик сообщений WM_COMMAND, распределяющий логику      *
* в зависимости от нажатого элемента интерфейса.                         *
*************************************************************************/
void HandleWmCommand(HWND hWnd, WPARAM wParam) {
    int wmId = LOWORD(wParam);
    if (wmId == ID_BTN_IMPORT || wmId == ID_BTN_SHOW) 
        HandleCmdImportShow(hWnd, wmId);
    else if (wmId == ID_BTN_CALC || wmId == ID_BTN_EXPORT) 
        HandleCmdCalcExport(hWnd, wmId);
    else if (wmId == ID_BTN_VIS_A || wmId == ID_BTN_VIS_B) 
        HandleCmdVis(hWnd, wmId);
    InvalidateRect(hWnd, NULL, TRUE);
}
\end{listing}
\normalsize

\pagebreak

\section{Функция WndProc}
\footnotesize
\begin{listing}{1}
/*************************************************************************
* Функция       : WndProc                                                *
*------------------------------------------------------------------------*
* Назначение    : Основная процедура обработки событий окна              *
* Входные       : hWnd, msg, wParam, lParam (стандартные Win32)          *
* Выходные      : LRESULT - результат обработки события                  *
*------------------------------------------------------------------------*
* ОПИСАНИЕ ФУНКЦИОНИРОВАНИЯ:                                             *
* Главный цикл сообщений приложения, управляет жизненным циклом окна,    *
* отрисовкой и взаимодействием с пользователем.                          *
*************************************************************************/
LRESULT CALLBACK WndProc(HWND hWnd, UINT msg, WPARAM wParam, LPARAM lParam) {
    switch (msg) { case WM_CREATE: 
        CreateInterfaceButtons(hWnd, ((LPCREATESTRUCT)lParam)->hInstance); 
		break;
    case WM_COMMAND:  HandleWmCommand(hWnd, wParam); 
		break;
    case WM_PAINT: HandlePaint(hWnd); 
		break;
    case WM_DESTROY: if (g_hPlotBitmapA) DeleteObject(g_hPlotBitmapA);
        if (g_hPlotBitmapB) DeleteObject(g_hPlotBitmapB); PostQuitMessage(0); 
		break;
    default:  return DefWindowProc(hWnd, msg, wParam, lParam);}
    return 0;
}
\end{listing}
\normalsize

\pagebreak%
\section{Функция WinMain}
\footnotesize
\begin{listing}{1}
/*************************************************************************
* Функция       : WinMain                                                *
*------------------------------------------------------------------------*
* Назначение    : Точка входа в Windows-приложение                       *
* Входные       : hInstance, hPrevInstance, lpCmdLine, nCmdShow          *
* Выходные      : int - код завершения программы                         *
*------------------------------------------------------------------------*
* ОПИСАНИЕ ФУНКЦИОНИРОВАНИЯ:                                             *
* Регистрирует класс окна, создает основное окно приложения и запускает  *
* цикл обработки сообщений (Message Loop).                                *
*************************************************************************/
int WINAPI WinMain(_In_ HINSTANCE hInst, _In_opt_ HINSTANCE hPInst, 
                   _In_ LPSTR cmd, _In_ int nCmdShow) {
    WNDCLASSEXW wcex = { sizeof(WNDCLASSEX), CS_HREDRAW | CS_VREDRAW, 
		WndProc, 0, 0, hInst, NULL, LoadCursor(NULL, IDC_ARROW), 
		(HBRUSH)(COLOR_WINDOW + 1), NULL, L"FRC", NULL };
    if (!RegisterClassExW(&wcex)) return 0;
    HWND hWnd = CreateWindowW(L"FRC", 
		L"Теория Поля — Расчет параметров векторных полей", WS_OVERLAPPEDWINDOW,
		CW_USEDEFAULT, CW_USEDEFAULT, 1200, 800, NULL, NULL, hInst, NULL);                     
    ShowWindow(hWnd, nCmdShow); UpdateWindow(hWnd); MSG msg; 
    while (GetMessage(&msg, NULL, 0, 0)) { 
        TranslateMessage(&msg); DispatchMessage(&msg); 
    }
    return (int)msg.wParam;
}
\end{listing}
\normalsize

\pagebreak

\section{Функция PolyToWolframString}
\footnotesize
\begin{listing}{1}
/*************************************************************************
* Функция       : PolyToWolframString                                    *
*------------------------------------------------------------------------*
* Назначение    : Конвертация многочлена в формат Wolfram Language      *
* Входные       : p - объект Polynomial                                  *
* Выходные      : wstring - строка для Mathematica/WolframScript         *
*------------------------------------------------------------------------*
* ОПИСАНИЕ ФУНКЦИОНИРОВАНИЯ:                                             *
* Преобразует внутреннее представление полинома в математическую строку, *
* понятную для интерпретатора Wolfram.                                   *
*************************************************************************/
wstring PolyToWolframString(const Polynomial& p) {
    if (p.terms.empty()) return L"0"; wstringstream ss; bool first = true;
    for (const auto& term : p.terms) { if (term.coef > 0 && !first) ss << L" + ";
        if (term.coef < 0) ss << (first ? L"-" : L" - ");
        ss << fixed << setprecision(4) << abs(term.coef);
        if (term.px > 0) { 
            ss << L"*x"; if (term.px > 1) ss << L"^" << term.px;}
        if (term.py > 0) { 
            ss << L"*y"; if (term.py > 1) ss << L"^" << term.py;}
        if (term.pz > 0) { 
            ss << L"*z"; if (term.pz > 1) ss << L"^" << term.pz; 
        }
        first = false;
    }
    return ss.str();
}
\end{listing}
\normalsize

\pagebreak

\section{Функция GenerateWolframPlot}\footnotesize
\begin{listing}{1}
/*************************************************************************
* Функция        : GenerateWolframPlot                                   *
*------------------------------------------------------------------------*
* Назначение     : Координация процесса генерации и загрузки графика     *
* Входные        : isA - флаг выбора поля (true - поле A, false - поле B)*
*                : err - ссылка для записи кода ошибки в случае неудачи  *
* Выходные       : bool - true при успехе, false при возникновении ошибки*
*------------------------------------------------------------------------*
* ОПИСАНИЕ ФУНКЦИОНИРОВАНИЯ:                                             *
* Получает строковое представление компонентов поля, определяет имена    *
* файлов скрипта и графики, вызывает функцию записи скрипта WriteScript. *
* Затем через RunWolfram запускает рендеринг и загружает результирующий  *
* BMP-файл в соответствующий глобальный дескриптор.                      *
*************************************************************************/
bool GenerateWolframPlot(bool isA, ErrCode& err) {
    auto& fmc = FieldMathCalculations;
    wstring px = PolyToWolframString(isA ? fmc::fieldA_X : fmc::GetFieldBX());
    wstring py = PolyToWolframString(isA ? fmc::fieldA_Y : fmc::GetFieldBY());
    wstring pz = PolyToWolframString(isA ? fmc::fieldA_Z : fmc::GetFieldBZ());
    
    wstring scr = isA ? L"t_a.wl" : L"t_b.wl";
    wstring bmp = isA ? L"t_a.bmp" : L"t_b.bmp";
    DeleteFileW(bmp.c_str());
    
    if (!WriteScript(isA, px, py, pz, scr, bmp)) {
        err = ErrWlScript;
        return false;
    }
    if (!RunWolfram(scr)) {
        err = ErrWlNotFound;
        return false;
    }
    
    HBITMAP hBmp = (HBITMAP)LoadImageW(
        NULL, bmp.c_str(), IMAGE_BITMAP, 0, 0, LR_LOADFROMFILE
    );
    if (!hBmp) {
        err = ErrWlImage;
        return false;
    }
    
    if (isA) g_hPlotBitmapA = hBmp; else g_hPlotBitmapB = hBmp;
    return true;
}
\end{listing}
\normalsize

\pagebreak

\section{Функция WriteScript}
\footnotesize
\begin{listing}{1}
/*************************************************************************
* Функция        : WriteScript                                           *
*------------------------------------------------------------------------*
* Назначение     : Создание файла скрипта WolframScript (.wl)            *
* Входные        : isA - флаг выбора поля (true - поле A, false - поле B)*
*                : px, py, pz - компоненты вектора в формате Wolfram     *
*                : scr - путь для сохранения файла скрипта               *
*                : bmp - целевой путь сохранения результирующего графика *
* Выходные       : bool - true, если файл успешно записан, иначе false   *
*------------------------------------------------------------------------*
* ОПИСАНИЕ ФУНКЦИОНИРОВАНИЯ:                                             *
* Рассчитывает потенциал поля, генерирует код на языке Wolfram Language  *
* для построения комбинированного графика VectorPlot3D и ContourPlot3D,  *
* перекодирует полученную строку в UTF-8 и записывает её в файл скрипта. *
*************************************************************************/
bool WriteScript(bool isA, const wstring& px, const wstring& py,
                 const wstring& pz, const wstring& scr, const wstring& bmp) {
    wstring pot = isA ? L"Norm[f]" :
        PolyToWolframString(FieldMathCalculations::GetAnalyticalPotential());
    
    wstringstream ss;
    ss << L"f={" << px << L"," << py << L"," << pz << L"};\npot=" << pot << L";\n"
       << L"v=VectorPlot3D[f,{x,-2,2},{y,-2,2},{z,-2,2},VectorPoints->{7,7,7},"
       << L"VectorColorFunction->\"Rainbow\",VectorStyle->\"Arrow\",VectorSizes->0.7,"
       << L"VectorAspectRatio->0.1,ViewPoint->{2.5,-2.5,1.8}];\n"
       << L"s=ContourPlot3D[pot,{x,-2,2},{y,-2,2},{z,-2,2},Contours->1,"
       << L"ColorFunction->\"Rainbow\",ContourStyle->Directive[Opacity[0.7]],"
       << L"Mesh->None,ViewPoint->{2.5,-2.5,1.8}];\n"
       << L"combined=GraphicsRow[{s,v},ImageSize->1600];\n"
       << L"Export[\"" << bmp << L"\",combined,ImageResolution->150];\n";
    
    wstring ws = ss.str();
    int len = WideCharToMultiByte(CP_UTF8, 0, ws.c_str(), -1, 0, 0, 0, 0);
    string u8(len, '\0');
    WideCharToMultiByte(CP_UTF8, 0, ws.c_str(), -1, &u8[0], len, 0, 0);
    
    ofstream f(scr, ios::binary);
    if (!f.is_open()) return false;
    f.write(u8.c_str(), len - 1);
    return true;
}
\end{listing}
\normalsize

\pagebreak

\section{Функция RunWolfram}
\footnotesize
\begin{listing}{1}
/*************************************************************************
* Функция        : RunWolfram                                            *
*------------------------------------------------------------------------*
* Назначение     : Исполнение сгенерированного .wl скрипта во внешней    *
*                : среде WolframScript                                   *
* Входные        : scr - путь к запускаемому файлу скрипта               *
* Выходные       : bool - true, если процесс завершился успешно, иначе   *
*                : false                                                 *
*------------------------------------------------------------------------*
* ОПИСАНИЕ ФУНКЦИОНИРОВАНИЯ:                                             *
* Формирует команду вызова исполняемого файла WolframScript.exe с флагом *
* -file. Запускает скрытый фоновый процесс без создания окна консоли,    *
* ожидает окончания рендеринга графики и корректно закрывает дескрипторы *
*************************************************************************/
bool RunWolfram(const wstring& scr) {
    wstring cmd = L"\"C:\\Program Files\\Wolfram Research"
        L"\\WolframScript\\wolframscript.exe\" -file " + scr;
    
    STARTUPINFOW si = { sizeof(si) };
    si.dwFlags = STARTF_USESHOWWINDOW;
    si.wShowWindow = SW_HIDE;
    
    PROCESS_INFORMATION pi = { 0 };
    if (!CreateProcessW(
        NULL, &cmd[0], NULL, NULL, FALSE,
        CREATE_NO_WINDOW, NULL, NULL, &si, &pi
    )) return false;
    
    WaitForSingleObject(pi.hProcess, INFINITE);
    CloseHandle(pi.hProcess);
    CloseHandle(pi.hThread);
    return true;
}
\end{listing}
\normalsize

\pagebreak

\section{Процедура DrawVisualization}\footnotesize
\begin{listing}{1}
/*************************************************************************
* Процедура      : DrawVisualization                                     *
*------------------------------------------------------------------------*
* Назначение     : Отрисовка графического представления выбранного поля   *
* Входные        : hdc  - дескриптор контекста устройства вывода         *
*                : rect - прямоугольная область для рисования            *
* Выходные       : Нет                                                   *
*------------------------------------------------------------------------*
* ОПИСАНИЕ ФУНКЦИОНИРОВАНИЯ:                                             *
* Определяет активную битовую карту поля, рассчитывает пропорциональные *
* координаты вывода с помощью GetDrawCoords, масштабирует изображение   *
* (StretchBlt с качественным сглаживанием) и выводит подписи к графикам.*
*************************************************************************/
void DrawVisualization(HDC hdc, RECT rect) {
    HBITMAP hBmp = (g_ActiveView == 3) ? g_hPlotBitmapA : g_hPlotBitmapB;
    if (!hBmp) return;
    
    HDC hdcMem = CreateCompatibleDC(hdc);
    HBITMAP hOldBmp = (HBITMAP)SelectObject(hdcMem, hBmp);
    
    BITMAP bitmap;
    GetObject(hBmp, sizeof(BITMAP), &bitmap);
    
    int lblH = 35, dW, dH, dX, dY;
    GetDrawCoords(bitmap, rect, lblH, dW, dH, dX, dY);
    
    SetStretchBltMode(hdc, HALFTONE);
    StretchBlt(
        hdc, dX, dY, dW, dH, hdcMem, 0, 0,
        bitmap.bmWidth, bitmap.bmHeight, SRCCOPY
    );
    
    SelectObject(hdcMem, hOldBmp);
    DeleteDC(hdcMem);
    
    DrawLabels(hdc, dX, dW, dY, dH, lblH);
}
\end{listing}
\normalsize

\pagebreak

\section{Процедура GetDrawCoords}
\footnotesize
\begin{listing}{1}
/*************************************************************************
* Процедура      : GetDrawCoords                                         *
*------------------------------------------------------------------------*
* Назначение     : Расчет координат вывода с сохранением пропорций       *
* Входные        : bmp  - структура с характеристиками изображения       *
*                : r    - прямоугольник доступной области рисования      *
*                : lblH - высота области, резервируемой под подписи      *
* Выходные       : dW, dH - рассчитанные ширина и высота для вывода     *
*                : dX, dY - рассчитанные координаты левого верхнего угла *
*------------------------------------------------------------------------*
* ОПИСАНИЕ ФУНКЦИОНИРОВАНИЯ:                                             *
* Сравнивает соотношение сторон исходного изображения и области вывода  *
* (с вычетом места под текст). Центрирует изображение по горизонтали     *
* или вертикали в зависимости от того, по какой оси произошло сжатие.   *
*************************************************************************/
void GetDrawCoords(const BITMAP& bmp, const RECT& r, int lblH,
                   int& dW, int& dH, int& dX, int& dY) {
    int rW = r.right - r.left;
    int rH = r.bottom - r.top - lblH;
    if (rH < 50) rH = 50;
    
    double aspB = (double)bmp.bmWidth / bmp.bmHeight;
    double aspR = (double)rW / rH;
    
    if (aspB > aspR) {
        dW = rW;
        dH = (int)(rW / aspB);
        dX = r.left;
        dY = r.top + (rH - dH) / 2;
    } else {
        dW = (int)(rH * aspB);
        dH = rH;
        dX = r.left + (rW - dW) / 2;
        dY = r.top;
    }
}
\end{listing}
\normalsize

\pagebreak

\section{Процедура DrawLabels}
\footnotesize
\begin{listing}{1}
/*************************************************************************
* Процедура      : DrawLabels                                            *
*------------------------------------------------------------------------*
* Назначение     : Отрисовка текстовых подписей под графиками поля       *
* Входные        : hdc  - контекст устройства для рисования текста       *
*                : dX, dW, dY, dH - геометрия выведенного графика        *
*                : lblH - высота, выделенная под отрисовку текста        *
* Выходные       : Нет                                                   *
*------------------------------------------------------------------------*
* ОПИСАНИЕ ФУНКЦИОНИРОВАНИЯ:                                             *
* Инициализирует шрифт Segoe UI Bold, настраивает прозрачный фон текста  *
* и выводит симметричные подписи (для эквипотенциалей слева и для       *
* силовых линий справа) под отмасштабированным изображением.            *
*************************************************************************/
void DrawLabels(HDC hdc, int dX, int dW, int dY, int dH, int lblH) {
    bool isA = (g_ActiveView == 3);
    wstring lL = isA ? L"Поле A: эквипотенциальные поверхности" 
                     : L"Поле B: эквипотенциальные поверхности";
    wstring lR = isA ? L"Поле A: силовые линии" : L"Поле B: силовые линии";
    
    HFONT hF = CreateFontW(
        16, 0, 0, 0, FW_BOLD, 0, 0, 0, DEFAULT_CHARSET,
        0, 0, CLEARTYPE_QUALITY, 0, L"Segoe UI"
    );
    HFONT hOldF = (HFONT)SelectObject(hdc, hF);
    
    int oldBk = SetBkMode(hdc, TRANSPARENT);
    COLORREF oldCol = SetTextColor(hdc, RGB(50, 50, 50));
    int lY = dY + dH + 5;
    
    RECT rL = { dX, lY, dX + dW / 2, lY + lblH };
    RECT rR = { dX + dW / 2, lY, dX + dW, lY + lblH };
    
    DrawTextW(hdc, lL.c_str(), -1, &rL, DT_CENTER | DT_TOP | DT_SINGLELINE);
    DrawTextW(hdc, lR.c_str(), -1, &rR, DT_CENTER | DT_TOP | DT_SINGLELINE);
    
    SetTextColor(hdc, oldCol);
    SetBkMode(hdc, oldBk);
    SelectObject(hdc, hOldF);
    DeleteObject(hF);
}
\end{listing}
\normalsize

\pagebreak

\section{Процедура ShowError}
\footnotesize
\begin{listing}{1}
/*************************************************************************
* Процедура     : ShowError                                              *
*------------------------------------------------------------------------*
* Назначение    : Централизованный вывод сообщений об ошибках            *
* Входные       : hWnd - окно, code - код ошибки                         *
* Выходные      : Нет                                                    *
*------------------------------------------------------------------------*
* ОПИСАНИЕ ФУНКЦИОНИРОВАНИЯ:                                             *
* Отображает MessageBox с соответствующим заголовком и текстом ошибки    *
* на основе переданного кода ошибки (enum).                              *
*************************************************************************/
void ShowError(HWND hWnd, ErrCode code) {
    const wchar_t* m = L"Неизвестная ошибка"; const wchar_t* t = L"Ошибка";
	UINT ic = MB_ICONERROR;
    if (code == ErrParse) { m = L"Файл поврежден, пуст или имеет неверный формат!"; 
        t = L"Ошибка парсинга"; }
    else if (code == WarnNoData) { m = L"Сперва необходимо загрузить данные!"; 
        t = L"Предупреждение"; ic = MB_ICONWARNING; }
    else if (code == WarnNoCalc) { m = L"Сперва необходимо провести расчеты!"; 
        t = L"Предупреждение"; ic = MB_ICONWARNING; }
    else if (code == ErrWrite) { m = L"Ошибка записи файла!"; }
    else if (code == ErrWlScript) { 
       m = L"Не удалось создать файл скрипта (.wl)! Проверьте права."; 
        t = L"Ошибка визуализации"; }
    else if (code == ErrWlNotFound) { 
       m = L"WolframScript не найден!\nПроверьте путь к wolframscript.exe."; 
        t = L"Ошибка визуализации"; }
    else if (code == ErrWlImage) { 
       m = L"Изображение не создано!\nВозможно, скрипт завершился с ошибкой."; 
        t = L"Ошибка визуализации"; }
MessageBoxW(hWnd, m, t, ic | MB_OK);
}
\end{listing}
\normalsize
\pagebreak%

\pagebreak

\label{lastpage}

\addcontentsline{toc}{chapter}{Лист регистрации изменений}
\renewcommand{\headrulewidth}{0pt}
\fontsize{3.5mm}{3mm}\selectfont
\noindent\begin{tblr}{colspec={|[\rs]p{8mm,c}|X[20,c]|X[20,c]|X[20,c]|X[20,c]|X[20,c]|X[25,c]|X[25,c]|X[15,c]|X[12,c]|[2pt]},
             rows   = {5mm},
             row{1} = {10mm,m},
             row{2} = {5mm,m},
             row{3} = {20mm,m},
             rowsep = 0pt,colsep = 0pt,
             vline{2-10} = {2-3}{\rs}
}\hline[\rs]
\SetCell[c=10]{c} {\fontsize{7mm}{0mm}\selectfont Лист регистрации изменени\U{й}} &&&&&&&&&&& \\\hline[\rs]
\SetCell[c=5]{c} {\small Номера листов (страниц)} &&&&& \SetCell[r=2]{c} Всего листов (страниц) в документе & \SetCell[r=2]{c} № документа & \SetCell[r=2]{c} {Входящи\U{й} № сопроводи-\\тельного  документа и дата} & \SetCell[r=2]{c} Подп. & \SetCell[r=2]{c} Дата \\\hline[\rs]
Изм & {Изменен-\\ных} & {Заменен-\\ных} & Новых & {Анулиро-\\ванных} &&&&& \\\hline[\rs]
&&&&&&&&&&& \\ \hline
&&&&&&&&&&& \\ \hline
&&&&&&&&&&& \\ \hline
&&&&&&&&&&& \\ \hline
&&&&&&&&&&& \\ \hline
&&&&&&&&&&& \\ \hline
&&&&&&&&&&& \\ \hline
&&&&&&&&&&& \\ \hline
&&&&&&&&&&& \\ \hline
&&&&&&&&&&& \\ \hline
&&&&&&&&&&& \\ \hline
&&&&&&&&&&& \\ \hline
&&&&&&&&&&& \\ \hline
&&&&&&&&&&& \\ \hline
&&&&&&&&&&& \\ \hline
&&&&&&&&&&& \\ \hline
&&&&&&&&&&& \\ \hline
&&&&&&&&&&& \\ \hline
&&&&&&&&&&& \\ \hline
&&&&&&&&&&& \\ \hline
&&&&&&&&&&& \\ \hline
&&&&&&&&&&& \\ \hline
&&&&&&&&&&& \\ \hline
&&&&&&&&&&& \\ \hline
&&&&&&&&&&& \\ \hline
&&&&&&&&&&& \\ \hline
&&&&&&&&&&& \\ \hline
&&&&&&&&&&& \\ \hline
&&&&&&&&&&& \\ \hline
&&&&&&&&&&& \\ \hline
&&&&&&&&&&& \\ \hline
&&&&&&&&&&& \\ \hline
&&&&&&&&&&& \\ \hline
&&&&&&&&&&& \\ \hline
&&&&&&&&&&& \\ \hline
&&&&&&&&&&& \\ \hline
&&&&&&&&&&& \\ \hline
&&&&&&&&&&& \\ \hline
&&&&&&&&&&& \\ \hline
&&&&&&&&&&& \\ \hline[\rs]
\end{tblr}

\end{document}